% octave_renorm.tex -- round 468, octave renormalization.
% OCTAVE_NOT_CAUCHY / EDGE_TWO_POINTS / REDUCTION_ALIAS.
% Builds with pdflatex.
% Companion machine check: verify_octave_renorm.py.
% Discovery: experiments/tfpt-discovery/octave_renorm_probe.py.
% NO RH CLAIM.  NO anti-RH claim.
\documentclass[11pt]{article}

\usepackage[a4paper,margin=2.7cm]{geometry}
\usepackage{amsmath,amssymb,amsthm}
\usepackage{microtype}
\usepackage{booktabs}
\usepackage{xcolor}
\usepackage[colorlinks=true,linkcolor=blue!50!black,urlcolor=blue!50!black]{hyperref}

\newtheorem{lemma}{Lemma}
\newtheorem{prop}{Proposition}

\title{\bfseries Octave renormalization\\[2mm]
\large intra-block $q^{\dagger}$ is not a Cauchy sequence}
\author{Stefan Hamann}
\date{August 31, 2026}

\begin{document}
\maketitle

\begin{abstract}\noindent
The named selected family has constant $\Delta$ on each
$r=\lfloor\sqrt k\rfloor$ block.
The live arithmetic source is $n<a_k$ (tents vanish for
$n\ge a_k$, r465), not the $a^2$-octave
$(a_k^2,a_{k+1}^2]$.
The active-octave weight $W$ and its unconditional
Chebyshev bound both \emph{grow}.
Isolating the new primes at fixed geometry can move
$q^{\dagger}$ by $0.22$ while the total step is
$-0.012$.
Intra-block $\Delta q$ oscillates through $k=15$.
Two $\Delta$-halving edges are positive
($+0.045$, $+0.018$) --- not a law.
Without a proved summable rate the proposed cut
``$q(\mathrm{block\text{-}limit}\,r)<1$ for infinitely
many $r$'' is an alias of the original quantifier.
No RH claim.
\end{abstract}

\medskip\noindent
\textbf{Status.}\enspace
No new SATZ.
Ausgang
\texttt{OCTAVE\_NOT\_CAUCHY / EDGE\_TWO\_POINTS / REDUCTION\_ALIAS}.
No $L^{*}$ claim.  No $R^{\dagger}$ claim.
No RH claim.  No anti-RH claim.

\section{Geometry}\label{sec:geom}

\begin{lemma}[$r$-sequence]\label{lem:r}
$r_k=\lfloor\sqrt k\rfloor$,
$\Delta_k=2^{-r_k}\log 2$,
$m_k=k\,2^{r_k}-1$.
In the measured range:
$k=5..8$ have $r=2$;
$k=9..15$ have $r=3$;
$k=16$ has $r=4$.
Block edges sit at $k=(r+1)^2$
($8\to9$, $15\to16$).
The step $9\to10$ is intra-block
($r=3\to3$), not an $r$-jump.
\end{lemma}

\begin{lemma}[active source]\label{lem:active}
On this mesh $M\Delta=\log a$, so every tent with
$\log n\ge\log a$ vanishes on lags $0,\ldots,M-1$
(r465).
The live increment $k\to k+1$ is the octave
$n\in[2^k,2^{k+1})$, plus $2^{r-1}$ new degrees
and a longer Archimedean grid.
The $a^2$-tail is a completeness label, not a
perturbation of $c^P$.
\end{lemma}

\section{Octave perturbation}\label{sec:A}

Unconditional Rosser--Schoenfeld:
$\psi(x)<1.03883\,x$.
Then
\[
\sum_{2^k\le n<2^{k+1}}\frac{\Lambda(n)}{\sqrt n}
\le
\frac{\psi_{\mathrm{up}}(2^{k+1})-\psi_{\mathrm{lo}}(2^k)}{2^{k/2}},
\]
which \emph{grows} in $k$.

\begin{center}
\scriptsize
\begin{tabular}{@{}rrrrrrc@{}}
\toprule
$k\to k{+}1$ & $\Delta q$ & $\Delta q_{\mathrm{oct}}$ &
 $W$ & bound & $dN$ & type \\
\midrule
$5\to6$  & $+0.00746$ & $+0.0669$ & $4.31$  & $11.8$  & $2$  & intra \\
$6\to7$  & $-0.02486$ & $+0.1437$ & $6.70$  & $10.5$  & $2$  & intra \\
$7\to8$  & $-0.00833$ & $+0.0417$ & $9.15$  & $14.5$  & $2$  & intra \\
$8\to9$  & $+0.04542$ & $+0.0989$ & $13.65$ & $20.1$  & $20$ & BLOCK \\
$9\to10$ & $+0.03439$ & $+0.0464$ & $18.63$ & $28.0$  & $4$  & intra \\
$10\to11$& $-0.01180$ & $+0.2199$ & $26.31$ & $39.1$  & $4$  & intra \\
$11\to12$& $-0.00624$ & $+0.0994$ & $37.81$ & $54.7$  & $4$  & intra \\
$12\to13$& $+0.01185$ & $+0.0165$ & $52.68$ & $76.7$  & $4$  & intra \\
$13\to14$& $+0.01784$ & $+0.0323$ & $75.34$ & $107.6$ & $4$  & intra \\
$14\to15$& $-0.01703$ & $+0.0068$ & $105.7$ & $151.1$ & $4$  & intra \\
$15\to16$& $+0.01844$ & $+0.0091$ & $149.9$ & $212.5$ & $68$ & BLOCK \\
\bottomrule
\end{tabular}
\end{center}

$\Delta q_{\mathrm{oct}}$ is the new-octave effect at
\emph{fixed} $k{+}1$ geometry.
Sharpness $\mathrm{bound}/|\Delta q|$ lies in
$400$--$11\,000$.
Cauchy interlacing on a $2^{r-1}$-row pad does not
produce a contraction of $q^{\dagger}=s_N/(s_{N-1}+5/7)$.

\begin{lemma}[not Cauchy]\label{lem:nocauchy}
The octave-weight bound is unconditional and not
summable.
Measured intra-block $\Delta q$ oscillates in
$[-0.025,+0.034]$ through $k=15$ and is not
dominated by $W$.
\texttt{OCTAVE\_NOT\_CAUCHY}.
\end{lemma}

\section{Block edges}\label{sec:B}

\begin{lemma}[two points]\label{lem:edge}
The only $\Delta$-halvings in range are
$8\to9$ ($+0.04542$) and $15\to16$ ($+0.01844$).
Both signs are $+$.
The ratio $2.46$ is not $2$.
\texttt{EDGE\_TWO\_POINTS}.
\end{lemma}

New indices carrying the table were interval-certified
by the r465 Chebyshev--Cholesky route
($k=13,14,15,16$ all have $q_{\mathrm{hi}}<1$;
$k=16$ took $35\,\mathrm{s}$).
$k=17..20$ were not needed for the verdict
(next edge is $k=25$).

\section{The reduced problem}\label{sec:C}

\paragraph{Wording.}
``$q^{\dagger}$ of the $r$-block limit is $<1$ for
infinitely many $r$, using a proved intra-block
Cauchy rate.''

\paragraph{What is proved.}
The $r$-sequence; vanishing of $n\ge a$ tents;
an unconditional growing bound on the active-octave
mass.

\paragraph{What is only measured.}
The $\Delta q$ table; two edge jumps; certified
enclosures at $k=13..16$.

\paragraph{Alias (r303).}
A finite block has $2r+1$ members.
Without a rate $\sum|\Delta q|<\infty$ the block
limit is not known to exist, and even if it existed
one would still need infinitely many living block
limits --- the same $\infty$ as infinitely many $k$,
up to a constant factor $2r+1$.
That is a slack shift, not a smaller quantifier.
\texttt{REDUCTION\_ALIAS}.

\section{Kill table}\label{sec:kills}

\begin{center}
\begin{tabular}{@{}lp{7.8cm}@{}}
\toprule
Kill & Status \\
\midrule
Comb inventory $n<a$ (live) vs $n\le a^2$ (label)
  & sealed; $k=5$ is $17$ vs $198$ \\
r465 certificates on new verdict $k$
  & $13,14,15,16$ live \\
$\psi$-bound vs measurement
  & bound grows; sharpness $400+$ \\
Determinism
  & two identical $k=5$ packs \\
$9\to10$ as $r$-jump
  & false; both $r=3$ \\
Quantifier cut
  & refused (alias) \\
\bottomrule
\end{tabular}
\end{center}

\medskip\noindent
\textbf{Claim boundary.}\enspace
Finite-window census.
No $L^{*}$ claim.
No $R^{\dagger}$ claim.
No RH claim.
No anti-RH claim.

\end{document}
